% SEGMENT OLP-0401-S01
% Part: many-valued-logic
% SOURCE CORRECTION TA-MVL-018: correct both the chapter and section metadata to infinite-valued-logics/goedel.
% Chapter: infinite-valued-logics
% Section: goedel

\documentclass[../../../include/open-logic-section]{subfiles}

\begin{document}

\olfileid{mvl}{inf}{god}

\olsection{கோடல்~(G\"odel) தருக்கங்கள்}

\begin{editorial}
  இது முடிவிலா மதிப்பு கோடல்~(G\"odel) தருக்கத்தைப் பற்றிய சுருக்கமான,
  முழுமையடையாத பிரிவு.
\end{editorial}

\begin{defn}\ollabel{def:goedel} முடிவிலா மதிப்பு கோடல்~(G\"odel)
தருக்கம்~$\LogGod[\infty]$ பின்வரும் அணியைக் கொண்டு வரையறுக்கப்படுகிறது:
\begin{enumerate}
  \item $\lfalse$, $\lnot$, $\land$, $\lor$, $\lif$ ஆகியவற்றைக் கொண்ட
  செந்தரக் கூற்று மொழி~$\Lang L_0$.
  \item மெய்மதிப்புகளின் கணம்~$V_\infty$.
  \item $1$ மட்டுமே நியமிக்கப்பட்ட மதிப்பு; அதாவது, $V^+ = \{1\}$.
  \item வாய்மைச் சார்புகள் பின்வரும் சார்புகளால் தரப்படுகின்றன:
% SOURCE CORRECTION TA-MVL-019: remove nested inline-math delimiters from the cases environment.
  \begin{align*}
    \tf{\lfalse} & = 0\\
    \tf{\lnot}[\LogGod](x) & = \begin{cases}
      1 & \text{எனில் } x =0\\
      0 & \text{வேறெனில்}
    \end{cases}\\
    \tf{\land}[\LogGod](x,y) & = \min(x,y)\\
    \tf{\lor}[\LogGod](x,y) & = \max(x,y)\\
    \tf{\lif}[\LogGod](x,y) & = \begin{cases}
      1 & \text{எனில் } x \le y\\
      y & \text{வேறெனில்.}
    \end{cases}
    \end{align*}
\end{enumerate}
$m$-மதிப்பு கோடல்~(G\"odel) தருக்கமும் இதேவாறு வரையறுக்கப்படுகிறது;
ஆனால் அதில் $V = V_m$.
\end{defn}

\begin{prop}
  \olref[thr][god]{defn:goedel} இல் வரையறுக்கப்பட்ட தருக்கம்
  $\LogGod[3]$, \olref{def:goedel} இல் வரையறுக்கப்பட்ட
  $\LogGod[3]$ உடன் ஒன்றே.
\end{prop}

\begin{proof}
  \olref[thr][god]{defn:goedel} இல் இணைப்பிகளுக்காகத் தரப்பட்ட வாய்மை
  அட்டவணைகளையும், \olref{def:goedel} இலுள்ள சமன்பாடுகள் தீர்மானிக்கும்
  பின்வரும் வாய்மை அட்டவணைகளையும் ஒப்பிட்டால் இதைக் காணலாம்:
  \begin{center}
    \begin{tabular}{c|c} 
      $\tf{\lnot}[\LogGod[3]]$ & \\ 
      \hline  
      $1$ & $0$ \\ 
      $1/2$ & $0$ \\
      $0$ & $1$ 
    \end{tabular}
    \quad
    \begin{tabular}{c|ccc} 
      $\tf{\land}[\LogGod]$ & $1$ & $1/2$ & $0$ \\ 
      \hline 
      $1$ & $1$ & $1/2$ & $0$ \\ 
      $1/2$ & $1/2$ & $1/2$ & $0$\\ 
      $0$ & $0$ & $0$ & $0$ 
    \end{tabular}
    \\[2ex]
    \begin{tabular}{c|ccc} 
      $\tf{\lor}[\LogGod]$ & $1$ & $1/2$ & $0$ \\ 
      \hline 
      $1$ & $1$ & $1$ & $1$ \\ 
      $1/2$ & $1$ & $1/2$ & $1/2$ \\
      $0$ & $1$ & $1/2$ & $0$ 
    \end{tabular}
    \quad
    \begin{tabular}{c|ccc} 
      $\tf{\lif}[\LogGod]$ & $1$ & $1/2$ & $0$ \\ 
      \hline 
      $1$ & $1$ & $1/2$ & $0$ \\ 
      $1/2$ & $1$ & $1$ & $0$  \\ 
      $0$ & $1$ & $1$ & $1$ 
    \end{tabular}
  \end{center} 
\end{proof}

\begin{prop}\ollabel{prop:god-infty-m}
  $\Gamma \Entails[\LogGod[\infty]] !B$ எனில், எல்லா~$m \ge 2$ க்கும்
  $\Gamma \Entails[\LogGod[m]] !B$.
\end{prop}

\begin{proof}
  பயிற்சி.
\end{proof}

\begin{prob}
  \olref[mvl][inf][god]{prop:god-infty-m} ஐ நிறுவுக.
\end{prob}

உண்மையில், இதன் மறுதலையும் அமைகிறது.

$\LogGod[3]$ ஐப் போலவே, $\LogGod[\infty]$ உம் உள்ளுணர்வுவாதத்தில்
செல்லுபடியாகும் எல்லா !!{formula}s ஐயும் மெய்மங்களாகக் கொண்டுள்ளது;
மேலும் பின்வரும் மெய்மமல்லாத எடுத்துக்காட்டுகள்~$\LogGod[\infty]$ இலும்
மெய்மங்கள் அல்ல:
\begin{align*}
  & p \lor \lnot p && (p \lif q) \lif (\lnot p \lor q) \\
  & \lnot\lnot p \lif p && \lnot(\lnot p \land \lnot q) \lif (p \lor q) \\
  & ((p \lif q) \lif p) \lif p && \lnot(p \lif q) \lif (p \land \lnot q)
\end{align*}
உள்ளுணர்வுவாதத்தில் செல்லுபடியாகாவிட்டாலும் $\LogGod[3]$ இன் மெய்மமாக
இருக்கும் $(p \lif q) \lor (q \lif p)$ என்ற !!{formula},
$\LogGod[\infty]$ இலும் ஒரு மெய்மம். உண்மையில், உள்ளுணர்வுவாதத்
தருக்கத்துடன் $(!A \lif !B) \lor (!B \lif !A)$ என்ற திட்டவடிவத்தைச்
சேர்த்த தருக்கமாக~$\LogGod[\infty]$ ஐ வகைப்படுத்தலாம். இதை மைக்கேல்
டம்மெட் நிறுவினார்; எனவே~$\LogGod[\infty]$ பெரும்பாலும் கோடல்--டம்மெட்
தருக்கம்~(G\"odel--Dummett logic)~$\Log{LC}$ எனக் குறிப்பிடப்படுகிறது.

\begin{prob}
  $(p \lif q) \lor (q \lif p)$ என்பது~$\LogGod[\infty]$ இன் ஒரு
  மெய்மம் என்று காட்டுக.
\end{prob}

\begin{prob}
  $\LogGod[3]$ இன் மெய்மமான
  $(p \lif q) \lor (q \lif r) \lor (r \lif s)$ என்பது
  $\LogGod[\infty]$ இன் மெய்மம் அல்ல என்று காட்டுக.
\end{prob}

\end{document}
